% !TeX root = ../papers/nc-p3.tex
\section{二十一次二重曲線からの根の一意復元と二次曲面射影の一般排除}
\label{sec:quadric-ruling-identifiability}
\begin{theorem}\label{thm:intrinsic-secant-scroll-root}
三次評価同型でどの直線にも四点がない被約二十点 $M$ を取り、三重消滅八次式の空間は一次元とする。
その八次曲面 $D$ は整、正規化は射影直線上の直線束の射影化（各繊維は入力直線に同型に写る）であるとし、
特異軌跡の唯一の曲線 $C$ は整な二十一次曲線で、一般横断平面の特異点は節点とする。
ここで直線繊維を持つ束は階数二ベクトル束の射影化を意味する。
このとき $M$ の十個の独立な二次式を成分に持ち、その二十点で平方となる根行列のうち、
一般階数四で $C$ に沿い階数二となるものが一つ存在すれば、出力の基底を除き非恒等根はそれだけである。
\end{theorem}
\begin{proof}
別の非恒等根 $N'$ を取る。定理\ref{thm:opposite-root-triple-octics}の被約点での階数排除は
三次評価の単射性・直線の印数・十成分の独立性で成立し、行列式は非零である。
従ってその行列式も $D$ を定める。階数二以下は $D$ の特異軌跡に入る。
それが有限なら定理\ref{thm:double-point-opposite-reconstruction}より先の根も同じVeronese像を持つ。
階数二曲線を持つ先の根と異なる有限型の根はこれでは生じない（同じ平方像なら出力合同である）。
正次元なら台は必ず $C$ である。一般点の階数は二、余因子スキームはそこで被約である。
実際階数一なら行列式の位数が三以上であり、階数二で二行消去後の三成分の一次部が
一方向しか張らなければ、二次初項が平方か零になり節点を与えない。

よって一般平面切断の核線束の次数は定理\ref{thm:cofactor-sectional-kernel}の局所計算により三である。
同じHodge指数計算と核像の非退化性は上記の仮定で使え、核像は有理正規三次曲線である。
一般平面切断上でこの曲線への射は次数一なので、一般核繊維の閉包は次数一の直線
（一般繊維は標数零で生成的被約、一平面との全交点数が一）である。
八次曲面の母線族は一意である。実際二つの異なる一径数の母線族があれば、
第一族の一般三本が互いにねじれの位置のとき第二族の全母線はこの三本に会い、滑らかな二次曲面に入る。
一般二本が常に会えば第一族は共面か共点であり、非平面の錐では頂点を通らない直線の存在が
その直線と頂点の張る平面を曲面の成分にしてしまう。
従って正規化 $\nu:S\to D$ 上で核点は既存の射影 $\pi:S\to\mathbb P^1$ と三次完全系で決まる。

$M'=\operatorname{coker}N'$ は $\nu_*\pi^*\mathcal O(3)$ に入る。四つの核座標が全域で射を与え、
$M'$ は一般階数一で捩れなしだからである。入力超平面を $H$ とすれば
$\pi_*\mathcal O(H)$ は次数八の階数二束（$H^2=8$）、$H$ の繊維次数は一である。
その $t$ 次対称冪の次数は $4t(t+1)$、階数は $t+1$ で、三次捻り後のEuler数は
$4(t+1)^2=\chi(M'(t))$。従って包含は等号である。同じ正規化上の線束の極小対称表示は
定理\ref{thm:quadric-kernel-projection}の自己準同型の議論で出力合同・スカラーを除き一意である。
一般階数一の根は恒等Veronese像であるため別の標識付き座標にはならない。
\end{proof}

\begin{theorem}\label{thm:generic-quadric-all-middles-exclusion}
滑らかな二次曲面の四空間 $V\in\operatorname{Gr}(4,S_{22})$ のある非空稠密開集合では、
二十個の相異なる長さ三繊維の入力点とそれらが張る反対平面を持つ
四辺有限標準正則窓は、中間座標の全ての選択について存在しない。
前節の有理標本はこの開集合に属する。
\end{theorem}
\begin{proof}[母線の割線条件を平面六次式で検査する]
前節の二方向の評価列による $T:S_t^4\to S_s^4$ を取る。$m$ 空間で
$T(t_0m)(z)=T(t_1m)(z)=0$ を $A(z)m=0$ と書く。$A$ は二行四列、$z$ の次数四である。
相異なる $z,w$ の母線が会う条件を
\[
 P(z+w,zw)=\det\begin{pmatrix}A(z)\\A(w)\end{pmatrix}/(z-w)^2
\]
と書く。対角で二行の差を取るので斉次表示では対角式の平方で割れ、
対称二次数 $(6,6)$ は $\operatorname{Sym}^2\mathbb P^1=\mathbb P^2$ の六次式に降りる。
$P$ が滑らか、対角二次曲線との十二交点が単純である条件を課す。
$A(z)$ の階数が下がれば $z$ を含む全二点因子が $P$ に入り、直線成分を生じてしまう。
異なる二本が同じ直線なら結合四行の階数は二以下で $P$ の特異点となる。
従って母線束 $S\to\mathbb P^1$ があり、入力への射は有限である。
実際固定 $m$ の逆像は、独立な二つの二項四次式の共通根なので有限である。
二本が会う組は曲線だから一般点で一本だけであり、$S$ は像の正規化となる。
Plücker式は基点を持たない八次式で、像は八次曲面である。それは二つの四次式の結果式、
従って $N_m=\operatorname{Bez}(T(t_0m),T(t_1m))$ の行列式である。
各辺の両射影が有限である開条件も課す。三本以上の母線が会う点は有限である。
共通三次因子は前節の標準 $m,y$ 辺の点を与えるからである。

対角の外の $P$ は異なる二本の交点を一意に決め、正規化の二点が同じ入力へ写る軌跡を
局所的に順序を除いて表示する。実際結合行列の階数三の小行列式でその点を解ける。
一般点で二枝の接平面は異なる。両微分が単射で接平面が一致すると、
二つの像の差の微分は階数二以下となり、$S\times_{\mathbb P^3}S$ の当該点の接次元は二以上、
滑らかな割線条件と矛盾する。微分が単射でない入力点は有限である。
局所的に直線を二空間のグラフとして動かすと、その一次変化は二行写像 $J$ であり、
母線上で潰れる方向があるのは $J$ の核点である。$P$ の対角制限は単元を除き $\det J$ で、
単純零点では $J$ は階数一で核点がただ一つとなる。
従って唯一の特異曲線は整な $P$ の像であり、一般平面切断は節点以外は滑らか。
実際一般平面は三本以上の母線の交点、非浸入点、二枝が横断しない有限集合を避ける。
一枝で浸入なら有限射の像はその点で滑らかである。平面八次の正規化は各母線との交点で
$\mathbb P^1$ なので節点は二十一個、特異曲線の次数は二十一である。
その一般点でBezout行列は階数二である（相異なる二共通根を持ち、第三の根はない）。
\end{proof}

\begin{proof}[全中間座標への適用と開条件の実在]
$x$ の二十個の長さ三繊維から前節の $m,y,n$ を得る。
母線制限が全二次式なら各母線の写像は二次Veronese埋込みであり、
長さ三繊維が一本の母線に入ることはない。
$S_{23},S_{32}$ は各長さ三繊維へ全射する。実際三次元射影空間の二次式は
全ての長さ三部分スキームを分離し、$Q$ では $S_{22}$ を与え、さらに消えない一次因子を掛けられる。
有限な四つの双一次辺は二十点の完全交叉だから、得た二十組のうち
$x,m$ のグラフでは他の点はなく、逆射影の繊維が線形なので $m$ も全て相異なる。
残りの辺でも同じく全ての組を尽くす。$m$ の二十点配置を $M$ と書けば、
その三次評価は同型であり、四点直線を持たない（有限辺の最大小行列式分解と直線切断上界）。
三重に通る八次式がスカラーを除いて一つという条件も課す。
$m,n$ の二つの三次因子は前節の $T$ により標準三割平面の反対対を与え、
$N_m$ は $M$ 上で $n$ の平方になる。十成分の独立性はBezoutの線形同型と
$\bigwedge^2 T$ の可逆性から従う。

任意の正則延長で $x$ の隣に現れる中間配置 $U$ を取る。
有限辺の完全交叉の跡による直交性から、その二次評価空間は標識点上の単元倍を除き $m^2$ に一致する。
従って $U$ は $M$ の全階数平方根であり、前定理から $m,n$ 以外の標識付き射影配置はない。
異なる二つを選べば前節の三割平面障害を受ける（反転にも左右正則性を適用）。
同じものを選べば中央の二辺は一辺とその転置に一致する。有限完全交叉を通る双一次式は
元の六式で尽きるからである。両方の候補でこの折り返しの三次交叉次元が四未満となる開条件を課せば、
これも標準窓に延長できない。

\path{computations/verify_quadric_ruling_identifiability.py} は有理標本で開条件を検算する。
母線方向の乗法の転置二本を、三次から四次への二つの二項乗法と同定して $T$ を作る。
良い素数1009で六次曲線の三偏導の八次倍が十三次全空間を張り、対角制限は十二次で平方自由。
$m$ の三重八次の二階微分評価は階数164で、結果式が標数零で存在するので核は一次元である。
同じ辺の三次・四次商でアフィン座標の特性多項式が平方自由で、二十印の被約性が従う。
両折り返しで交叉は有理数上零である。
また有理標本で $x,m$ 辺の最大小行列式と、対称射影根の二行小行列式全体とは
同じ十五次元の四次式空間を張ることを、三十五点での補間により有理数上で確かめる。
従ってここで入力の長さ三繊維は実際に完全交叉グラフの二十印を与える。
これらの検査に使った最大階数、滑らかさ、有限性の条件は開であり、
他の平方根を排除する議論自体は指定開集合内の各窓に適用できる。
\end{proof}
